\section{Related work}
\label{sec:related}
\subsection{LLM traveler models and empirical alignment}
LLMs have been used to generate activities and travel in urban simulation \citep{gatsim2026,genagents2026}, model day-to-day route choice \citep{wang2025agentic}, and predict modes using retrieved examples or incremental learning \citep{ragbenchmark2026,incremental2026}. Conceptual work has set out their possible roles within transport systems and motivated hybrid designs \citep{liu2025framework,nie2025roles}. These applications establish the relevance of LLM-generated decisions; they do not make fidelity to an LLM equivalent to fidelity to travelers.

Empirical alignment is a separate research problem. Studies of synthetic survey respondents document both reproducible aggregate patterns and sensitivity to prompts and model versions \citep{argyle2023out,aher2023using,bisbee2024synthetic}. For travel choice, \citet{valuingtime2026} examine the direction and magnitude of time trade-offs, while \citet{liu2026persona} learn persona-based representations using empirical choices. In this journal, \citet{xu2026satisfaction} compare zero-shot and few-shot LLM predictions of travel satisfaction and study improvements from task-specific examples. Learning and adjustment behavior have also been addressed explicitly \citep{dualagent2026}. These studies motivate empirical checks rather than treating the Teacher as a behavioral ground truth.

The distinction in the present study is the object being audited: a numerical Student fitted to archived LLM judgments. We evaluate whether its finite responses match those judgments and whether this ranking transfers to independently collected stated responses. The human comparison is an external evaluation, not a calibration procedure. The contemporary LLM subset is a later reference model and does not, without matched backend provenance, identify how much discrepancy originated during compression.

\subsection{Predictive and behavioral criteria for choice models}
Random-utility models connect choice probabilities to traveler characteristics and alternative attributes \citep{mcfadden1974conditional,ben1985discrete}. Generalized and mixed logit formulations represent richer heterogeneity \citep{walker2002generalized,hensher2003mixed,train2009discrete}, while representation learning can retain a specified utility core alongside flexible components \citep{sifringer2020enhancing}. Comparisons of these families show that predictive performance, market shares and behavioral indicators can favor different models \citep{wang2020deep,martinbaos2023prediction}.

Distillation changes the estimation target from observed choices to another model's elicited distributions. A structured utility model may preserve such responses through regularization and restricted functional form, without demonstrating that the Teacher or travelers follow that form. We therefore compare logit and neural specifications on the same information, then evaluate their human-response discrepancies separately. Independent timing predictors test whether choice and departure adjustment benefit from different specifications rather than assuming that a shared representation is preferable for both.

\subsection{Response supervision and simulation surrogates}
Knowledge distillation transfers predictions or representations to a smaller model \citep{hinton2015distilling,gou2021knowledge,sanh2019distilbert}. Relational distillation preserves relationships between examples \citep{park2019relational}; Sobolev training uses output and derivative information \citep{czarnecki2017sobolev}. Response-aware fitting is therefore not a new general principle. Here the relation is a finite, travel-specific difference between elicited distributions. It is calculated from endpoint targets already available to every controlled Student, so the experiment concerns an optimization preference for scenario contrasts rather than access to additional Teacher information.

Surrogate modeling also has an established transport role. \citet{natterer2025surrogates} approximate network outcomes of a MATSim model to accelerate scenario exploration. Our Student instead replaces the traveler decision function while retaining network construction and execution. This distinction makes assignment and routing part of the evaluation rather than outputs that the surrogate directly predicts. A concurrent preprint, CALM, combines calibrated choice, optional LLM planning, network feedback and replayable evaluation \citep{cheng2026calm}. Its system-level integration is complementary to our narrower question: whether finite responses survive compression and whether the resulting model ordering persists across Teacher, stated-choice and execution comparisons. Neither the use of an LLM in a simulator nor the presence of intervention tests alone constitutes the contribution claimed here.
